\documentclass[12pt, a4paper]{article}

% ── Packages ───────────────────────────────────────────────────────────
\usepackage[utf8]{inputenc}
\usepackage[T1]{fontenc}
\usepackage{lmodern}
\usepackage[margin=1in]{geometry}
\usepackage{hyperref}
\usepackage{xcolor}
\usepackage{listings}
\usepackage{graphicx}
\usepackage{booktabs}
\usepackage{enumitem}
\usepackage{parskip}
\usepackage{fancyhdr}
\usepackage{titlesec}

% ── Listings style ─────────────────────────────────────────────────────
\definecolor{codebg}{RGB}{245,245,250}
\definecolor{codeframe}{RGB}{180,180,200}
\definecolor{keyword}{RGB}{0,100,180}
\definecolor{string}{RGB}{170,55,55}
\definecolor{comment}{RGB}{100,140,100}

\lstset{
  backgroundcolor=\color{codebg},
  frame=single,
  rulecolor=\color{codeframe},
  basicstyle=\ttfamily\small,
  keywordstyle=\color{keyword}\bfseries,
  stringstyle=\color{string},
  commentstyle=\color{comment}\itshape,
  breaklines=true,
  breakatwhitespace=true,
  showstringspaces=false,
  tabsize=4,
  numbers=left,
  numberstyle=\tiny\color{gray},
  numbersep=8pt,
  xleftmargin=2em,
  framexleftmargin=1.5em,
  captionpos=b,
}

% ── Page style ─────────────────────────────────────────────────────────
\pagestyle{fancy}
\fancyhf{}
\rhead{\textit{AutoBrief Documentation}}
\lhead{\textit{\nouppercase{\leftmark}}}
\cfoot{\thepage}

\hypersetup{
  colorlinks=true,
  linkcolor=blue!70!black,
  urlcolor=blue!60!black,
  citecolor=green!50!black,
}

% ── Title ──────────────────────────────────────────────────────────────
\title{%
  \textbf{AutoBrief: Automated Daily Briefing Bot}\\[0.5em]
  \large Full-Stack Application --- Design, Architecture, and Implementation Guide
}
\author{AutoBrief Project Documentation}
\date{March 2026}

\begin{document}
\maketitle
\tableofcontents
\newpage

% ======================================================================
\section{Introduction}
% ======================================================================

AutoBrief is a full-stack web application that automates the generation and delivery of personalized daily briefings. It aggregates weather data, top news headlines, and a user's class schedule, then sends a formatted summary via Telegram at a configurable time each day.

\subsection{Key Design Principles}

\begin{itemize}[leftmargin=2em]
  \item \textbf{Zero-cost external APIs}: Weather via Open-Meteo (free, no API key), news via Google News RSS (free, no API key). Only Telegram requires a free bot token from BotFather.
  \item \textbf{Self-contained}: SQLite database requires no external database server. APScheduler runs inside the application process.
  \item \textbf{Preview without Telegram}: The application supports a preview mode --- users can see their briefing message without configuring Telegram.
\end{itemize}

% ======================================================================
\section{System Architecture}
% ======================================================================

The system follows a client-server architecture with three main layers:

\begin{enumerate}[leftmargin=2em]
  \item \textbf{Frontend} (React + Vite + Tailwind CSS): Single-page dashboard for configuration and manual triggering.
  \item \textbf{Backend} (Python + FastAPI): RESTful API server with CRUD endpoints, the briefing engine, and background scheduler.
  \item \textbf{Data Layer} (SQLite + SQLAlchemy): Persistent storage for user preferences and class schedule.
\end{enumerate}

\subsection{Architecture Diagram}

\begin{verbatim}
+-------------------+       HTTP/JSON        +-------------------+
|                   | <--------------------> |                   |
|   React Frontend  |   /api/* (proxied)     |  FastAPI Backend  |
|   (Vite :5173)    |                        |  (Uvicorn :8000)  |
|                   |                        |                   |
+-------------------+                        +--------+----------+
                                                      |
                                  +-------------------+-------------------+
                                  |                   |                   |
                           +------+------+    +-------+------+   +-------+-------+
                           |  SQLite DB  |    | Open-Meteo   |   | Google News   |
                           | (local file)|    | (free API)   |   | (RSS feed)    |
                           +-------------+    +--------------+   +---------------+
                                                      |
                                              +-------+-------+
                                              | Telegram Bot  |
                                              | API (free)    |
                                              +---------------+
\end{verbatim}

\subsection{Request Flow}

When a user clicks ``Preview Briefing'':

\begin{enumerate}[leftmargin=2em]
  \item React sends \texttt{POST /api/preview} to the backend.
  \item FastAPI instantiates a \texttt{BriefingEngine}.
  \item The engine reads user preferences from SQLite.
  \item It concurrently fetches weather (Open-Meteo) and news (Google News RSS).
  \item It reads today's class schedule from SQLite.
  \item It builds a Markdown message and returns it in the JSON response.
  \item React renders the message in a preview panel.
\end{enumerate}

% ======================================================================
\section{Directory Structure}
% ======================================================================

\begin{lstlisting}[language={}, title={Project Tree}]
AutoBrief/
  .env                        # Environment variables (Telegram token)
  .env.example                # Template for .env
  README.md                   # Quick-start guide
  autobrief_doc.tex           # This document
  backend/
    config.py                 # Loads env vars
    database.py               # SQLAlchemy engine, session, init
    models.py                 # ORM models (UserPreferences, ClassSchedule)
    briefing.py               # BriefingEngine (weather, news, Telegram)
    scheduler.py              # APScheduler cron job setup
    main.py                   # FastAPI app, routes, lifespan
    requirements.txt          # Python dependencies
    tests.py                  # Comprehensive test suite (23 tests)
    venv/                     # Python virtual environment
  frontend/
    index.html                # HTML entry point
    package.json              # Node.js dependencies
    vite.config.js            # Vite config with API proxy
    tailwind.config.js        # Tailwind CSS config
    postcss.config.js         # PostCSS config
    src/
      main.jsx                # React entry point
      index.css               # Tailwind directives
      App.jsx                 # Dashboard components
    node_modules/             # Node.js packages
\end{lstlisting}

% ======================================================================
\section{Technology Stack}
% ======================================================================

\begin{table}[h!]
\centering
\begin{tabular}{@{}lll@{}}
\toprule
\textbf{Layer} & \textbf{Technology} & \textbf{Purpose} \\
\midrule
Frontend       & React 18            & UI components and state management \\
               & Vite 6              & Build tool and dev server \\
               & Tailwind CSS 3      & Utility-first CSS framework \\
               & Axios               & HTTP client for API calls \\
\midrule
Backend        & Python 3.11+        & Server-side language \\
               & FastAPI 0.115       & Async web framework with auto docs \\
               & SQLAlchemy 2.0      & ORM for database access \\
               & APScheduler 3.10    & In-process cron-like scheduler \\
               & httpx 0.28          & Async HTTP client \\
               & feedparser 6.0      & RSS/Atom feed parser \\
\midrule
Database       & SQLite              & File-based relational database \\
\midrule
External APIs  & Open-Meteo          & Weather data (free, no key) \\
               & Google News RSS     & News headlines (free, no key) \\
               & Telegram Bot API    & Message delivery (free token) \\
\bottomrule
\end{tabular}
\caption{Technology stack overview}
\end{table}

% ======================================================================
\section{Database Models}
% ======================================================================

The application uses two SQLAlchemy ORM models backed by SQLite.

\subsection{UserPreferences}

Stores the user's configuration for the briefing engine. Only one row exists (singleton pattern --- upserted on save).

\begin{table}[h!]
\centering
\begin{tabular}{@{}llll@{}}
\toprule
\textbf{Column} & \textbf{Type} & \textbf{Nullable} & \textbf{Description} \\
\midrule
id               & Integer (PK)  & No  & Auto-increment primary key \\
city             & String        & No  & City name for weather lookup \\
news\_topics     & String        & No  & Comma-separated news search terms \\
telegram\_chat\_id & String      & No  & Telegram chat ID for delivery \\
\bottomrule
\end{tabular}
\caption{UserPreferences schema}
\end{table}

\begin{lstlisting}[language=Python, title={models.py --- UserPreferences}]
class UserPreferences(Base):
    __tablename__ = "user_preferences"

    id = Column(Integer, primary_key=True, index=True)
    city = Column(String, nullable=False, default="London")
    news_topics = Column(String, nullable=False, default="technology")
    telegram_chat_id = Column(String, nullable=False, default="")
\end{lstlisting}

\subsection{ClassSchedule}

Stores weekly class schedule entries. Multiple entries per day are supported, ordered by time.

\begin{table}[h!]
\centering
\begin{tabular}{@{}llll@{}}
\toprule
\textbf{Column} & \textbf{Type} & \textbf{Nullable} & \textbf{Description} \\
\midrule
id            & Integer (PK)  & No  & Auto-increment primary key \\
day\_of\_week & String        & No  & Day name (Monday--Sunday) \\
time          & String        & No  & 24h time string (e.g., ``09:00'') \\
subject       & String        & No  & Class/course name \\
location      & String        & No  & Room or building (optional) \\
\bottomrule
\end{tabular}
\caption{ClassSchedule schema}
\end{table}

% ======================================================================
\section{Backend Implementation}
% ======================================================================

\subsection{Configuration (\texttt{config.py})}

Environment variables are loaded from a \texttt{.env} file using \texttt{python-dotenv}. The only required secret is the Telegram bot token (and even that is optional for preview mode).

\begin{lstlisting}[language=Python, title={config.py}]
import os
from dotenv import load_dotenv

load_dotenv(os.path.join(os.path.dirname(__file__), "..", ".env"))

TELEGRAM_BOT_TOKEN = os.getenv("TELEGRAM_BOT_TOKEN", "")
BRIEFING_HOUR = int(os.getenv("BRIEFING_HOUR", "7"))
BRIEFING_MINUTE = int(os.getenv("BRIEFING_MINUTE", "0"))
DATABASE_URL = "sqlite:///" + os.path.join(
    os.path.dirname(__file__), "autobrief.db"
)
\end{lstlisting}

\subsection{Database Layer (\texttt{database.py})}

Uses SQLAlchemy 2.0 declarative base with a session factory. The \texttt{get\_db} generator provides dependency-injected sessions to FastAPI route handlers.

\begin{lstlisting}[language=Python, title={database.py}]
from sqlalchemy import create_engine
from sqlalchemy.orm import sessionmaker, DeclarativeBase
from config import DATABASE_URL

engine = create_engine(DATABASE_URL,
                       connect_args={"check_same_thread": False})
SessionLocal = sessionmaker(autocommit=False, autoflush=False,
                            bind=engine)

class Base(DeclarativeBase):
    pass

def get_db():
    db = SessionLocal()
    try:
        yield db
    finally:
        db.close()

def init_db():
    Base.metadata.create_all(bind=engine)
\end{lstlisting}

\subsection{Briefing Engine (\texttt{briefing.py})}

The core service class that orchestrates data fetching and message construction.

\subsubsection{Weather Fetcher --- Open-Meteo}

Uses a two-step process:
\begin{enumerate}[leftmargin=2em]
  \item \textbf{Geocoding}: Resolves city name to latitude/longitude via \texttt{geocoding-api.open-meteo.com}.
  \item \textbf{Forecast}: Fetches current weather using the coordinates via \texttt{api.open-meteo.com}.
\end{enumerate}

WMO weather interpretation codes are mapped to human-readable descriptions (e.g., code 3 becomes ``Overcast'').

\subsubsection{News Fetcher --- Google News RSS}

Fetches headlines by constructing a Google News RSS search URL with the user's topics, then parses the feed using the \texttt{feedparser} library. Article titles are split to extract the source name.

\subsubsection{Message Builder}

Constructs a Markdown-formatted message with three sections:
\begin{itemize}[leftmargin=2em]
  \item Weather: city, temperature, humidity, wind speed, conditions
  \item Headlines: numbered list of top 3 articles with sources
  \item Schedule: today's classes with times and locations
\end{itemize}

Each section gracefully handles missing data (e.g., ``unavailable'' for weather failures).

\subsubsection{Telegram Sender}

Posts the formatted message to the Telegram Bot API using \texttt{httpx}. Returns early with a descriptive result if the bot token is not configured.

\subsubsection{Orchestrator (\texttt{run} method)}

The \texttt{run(preview\_only=False)} method:
\begin{enumerate}[leftmargin=2em]
  \item Reads preferences from the database.
  \item Fetches weather and news (failures are caught individually).
  \item Reads today's class schedule.
  \item Builds the message.
  \item If \texttt{preview\_only=True} or no chat ID: returns the message without sending.
  \item Otherwise: sends via Telegram and returns the result.
\end{enumerate}

\subsection{Scheduler (\texttt{scheduler.py})}

Uses APScheduler's \texttt{BackgroundScheduler} with a \texttt{cron} trigger. The job runs daily at the configured hour and minute (default: 07:00).

Since APScheduler uses synchronous jobs but the \texttt{BriefingEngine} is async, a wrapper function creates a new event loop to run the coroutine:

\begin{lstlisting}[language=Python, title={scheduler.py}]
def _run_briefing_sync() -> None:
    engine = BriefingEngine()
    loop = asyncio.new_event_loop()
    try:
        result = loop.run_until_complete(engine.run())
        print(f"[Scheduler] Briefing result: {result}")
    finally:
        loop.close()

scheduler = BackgroundScheduler()
scheduler.add_job(
    _run_briefing_sync,
    trigger="cron",
    hour=BRIEFING_HOUR,
    minute=BRIEFING_MINUTE,
    id="daily_briefing",
    replace_existing=True,
)
\end{lstlisting}

\subsection{FastAPI Application (\texttt{main.py})}

\subsubsection{Lifespan Management}

The application uses FastAPI's async context manager lifespan to:
\begin{itemize}[leftmargin=2em]
  \item Initialize the database tables on startup.
  \item Start the APScheduler background scheduler.
  \item Shut down the scheduler on application exit.
\end{itemize}

\subsubsection{CORS Middleware}

Configured to allow requests from \texttt{http://localhost:5173} (the Vite dev server), with all methods and headers permitted.

\subsubsection{API Endpoints}

\begin{table}[h!]
\centering
\begin{tabular}{@{}llp{7cm}@{}}
\toprule
\textbf{Method} & \textbf{Path} & \textbf{Description} \\
\midrule
GET    & \texttt{/api/preferences}       & Returns saved preferences (or null) \\
POST   & \texttt{/api/preferences}       & Creates or updates the singleton preferences row \\
GET    & \texttt{/api/schedule}           & Lists all schedule entries; optional \texttt{?day=Monday} filter \\
POST   & \texttt{/api/schedule}           & Creates a new schedule entry \\
PUT    & \texttt{/api/schedule/\{id\}}    & Updates an existing schedule entry \\
DELETE & \texttt{/api/schedule/\{id\}}    & Deletes a schedule entry \\
POST   & \texttt{/api/trigger}            & Runs the full briefing pipeline (with Telegram send) \\
POST   & \texttt{/api/preview}            & Runs the briefing pipeline in preview mode (no send) \\
\bottomrule
\end{tabular}
\caption{REST API endpoints}
\end{table}

% ======================================================================
\section{Frontend Implementation}
% ======================================================================

The frontend is a single-page React application built with Vite and styled with Tailwind CSS.

\subsection{Development Proxy}

The Vite configuration proxies all \texttt{/api/*} requests to the FastAPI backend at \texttt{http://127.0.0.1:8000}, eliminating CORS issues during development:

\begin{lstlisting}[language=JavaScript, title={vite.config.js}]
export default defineConfig({
  plugins: [react()],
  server: {
    port: 5173,
    proxy: {
      "/api": "http://127.0.0.1:8000",
    },
  },
});
\end{lstlisting}

\subsection{Component Architecture}

The \texttt{App.jsx} file contains four components:

\begin{enumerate}[leftmargin=2em]
  \item \textbf{PreferencesForm}: Form with inputs for city, news topics, and Telegram chat ID. Loads existing preferences on mount, saves via POST.
  \item \textbf{ScheduleManager}: Combined form and table for CRUD operations on class schedule entries. Supports inline editing and deletion.
  \item \textbf{TriggerButton}: Two buttons --- ``Preview Briefing'' (calls \texttt{/api/preview}) and ``Send via Telegram'' (calls \texttt{/api/trigger}). Displays the formatted message in a preview panel.
  \item \textbf{App}: Shell component with header, three card-style sections, and footer.
\end{enumerate}

\subsection{State Management}

Each component manages its own local state using React's \texttt{useState} and \texttt{useEffect} hooks. API calls use Axios with a base URL of \texttt{/api} (proxied to the backend).

% ======================================================================
\section{External API Integration Details}
% ======================================================================

\subsection{Open-Meteo (Weather)}

\begin{itemize}[leftmargin=2em]
  \item \textbf{Cost}: Completely free, no API key required.
  \item \textbf{Geocoding endpoint}: \texttt{GET https://geocoding-api.open-meteo.com/v1/search?name=\{city\}\&count=1}
  \item \textbf{Forecast endpoint}: \texttt{GET https://api.open-meteo.com/v1/forecast} with parameters for latitude, longitude, and current weather.
  \item \textbf{Data returned}: Temperature (Celsius), wind speed (km/h), WMO weather code, hourly humidity.
  \item \textbf{Rate limits}: Fair use (no hard limit for reasonable usage).
\end{itemize}

\subsection{Google News RSS (News)}

\begin{itemize}[leftmargin=2em]
  \item \textbf{Cost}: Completely free, no API key required.
  \item \textbf{Endpoint}: \texttt{GET https://news.google.com/rss/search?q=\{query\}\&hl=en-US\&gl=US\&ceid=US:en}
  \item \textbf{Parser}: The \texttt{feedparser} library parses the RSS XML response.
  \item \textbf{Data returned}: Article title, source name (extracted from title suffix), and URL.
\end{itemize}

\subsection{Telegram Bot API (Delivery)}

\begin{itemize}[leftmargin=2em]
  \item \textbf{Cost}: Free. Requires creating a bot via BotFather (\texttt{@BotFather} on Telegram).
  \item \textbf{Setup}: Message BotFather with \texttt{/newbot}, follow prompts, receive a token.
  \item \textbf{Chat ID}: Start a conversation with the bot, then use \texttt{https://api.telegram.org/bot\{TOKEN\}/getUpdates} to find your chat ID.
  \item \textbf{Endpoint}: \texttt{POST https://api.telegram.org/bot\{TOKEN\}/sendMessage}
\end{itemize}

% ======================================================================
\section{Test Suite}
% ======================================================================

The project includes 23 automated tests in \texttt{tests.py}, organized into five test classes:

\begin{table}[h!]
\centering
\begin{tabular}{@{}llc@{}}
\toprule
\textbf{Test Class} & \textbf{Coverage} & \textbf{Count} \\
\midrule
TestModels             & SQLAlchemy model creation            & 2 \\
TestPreferencesAPI     & GET/POST preferences CRUD            & 3 \\
TestScheduleAPI        & Full CRUD + edge cases               & 7 \\
TestBriefingEngine     & Weather, news, message builder, run  & 8 \\
TestTriggerAPI         & Trigger and preview endpoints         & 3 \\
\midrule
\textbf{Total}         &                                      & \textbf{23} \\
\bottomrule
\end{tabular}
\caption{Test suite overview}
\end{table}

\subsection{Test Infrastructure}

\begin{itemize}[leftmargin=2em]
  \item Uses an in-memory SQLite database with \texttt{StaticPool} to share connections across threads.
  \item Each test gets fresh tables via an \texttt{autouse} fixture.
  \item FastAPI's dependency injection is overridden to use the test database.
  \item Integration tests hit the real Open-Meteo and Google News APIs (both free).
  \item \texttt{briefing.SessionLocal} is patched in orchestrator tests to use the test database.
\end{itemize}

Running the tests:

\begin{lstlisting}[language=bash, title={Running tests}]
cd backend
source venv/bin/activate
python -m pytest tests.py -v
\end{lstlisting}

% ======================================================================
\section{Setup and Execution Guide}
% ======================================================================

\subsection{Prerequisites}

\begin{itemize}[leftmargin=2em]
  \item Python 3.11 or newer
  \item Node.js 18 or newer
  \item (Optional) A Telegram bot token from BotFather (free)
\end{itemize}

\subsection{Environment Configuration}

Edit the \texttt{.env} file in the project root:

\begin{lstlisting}[language=bash, title={.env}]
TELEGRAM_BOT_TOKEN=your_telegram_bot_token
BRIEFING_HOUR=7
BRIEFING_MINUTE=0
\end{lstlisting}

The Telegram token is optional --- the app works in preview mode without it.

\subsection{Backend Setup}

\begin{lstlisting}[language=bash, title={Backend setup commands}]
cd backend
python3 -m venv venv
source venv/bin/activate        # macOS/Linux
pip install -r requirements.txt
uvicorn main:app --reload --port 8000
\end{lstlisting}

The API will be available at \texttt{http://127.0.0.1:8000}. Interactive API docs are at \texttt{http://127.0.0.1:8000/docs}.

\subsection{Frontend Setup}

\begin{lstlisting}[language=bash, title={Frontend setup commands}]
cd frontend
npm install
npm run dev
\end{lstlisting}

The dashboard will be available at \texttt{http://localhost:5173}.

\subsection{Usage Workflow}

\begin{enumerate}[leftmargin=2em]
  \item Open \texttt{http://localhost:5173} in a browser.
  \item Fill in preferences: city name, news topics, and (optionally) Telegram chat ID.
  \item Click ``Save Preferences''.
  \item Add class schedule entries using the schedule form.
  \item Click ``Preview Briefing'' to see the formatted message.
  \item (With Telegram configured) Click ``Send via Telegram'' to deliver it.
  \item The scheduler automatically sends the briefing daily at the configured time.
\end{enumerate}

% ======================================================================
\section{Sample Briefing Output}
% ======================================================================

Below is an actual briefing message generated by the system during testing:

\begin{lstlisting}[language={}, title={Sample briefing (March 16, 2026)}]
*Daily Briefing -- Monday, March 16, 2026*

*Weather*
City: New York
Temp: 15.0 C
Humidity: 69%
Wind: 26.2 km/h
Conditions: Overcast

*Top Headlines*
1. Zelenskiy says Ukraine wants money, technology
   in return for Middle East drone help (Reuters)
2. What Iran's war reveals about future of defense
   technology (The Jerusalem Post)
3. Gorilla Technology & Yotta Sign Landmark AI
   Infrastructure Deal (Gorilla Technology)

*Today's Classes (Monday)*
- 09:00  Data Structures -- Room 302

_Have a productive day!_
\end{lstlisting}

% ======================================================================
\section{Conclusion}
% ======================================================================

AutoBrief demonstrates a complete full-stack application with the following characteristics:

\begin{itemize}[leftmargin=2em]
  \item \textbf{Zero-cost operation}: All external APIs (Open-Meteo, Google News RSS) are completely free with no API keys required. Only Telegram needs a free bot token.
  \item \textbf{Self-contained deployment}: SQLite eliminates database server dependency. APScheduler runs in-process.
  \item \textbf{Comprehensive testing}: 23 automated tests covering models, API endpoints, the briefing engine, and integration with external APIs.
  \item \textbf{Graceful degradation}: Individual API failures are caught and the briefing still generates with available data.
  \item \textbf{Preview mode}: Full briefing generation without Telegram configuration for development and testing.
\end{itemize}

\end{document}
